\documentclass[../main]{subfiles}
\begin{document}

\chapter{Teorema de Bernoulli}

\section{Teoría del Teorema de Bernoulli}
\info{Teorema de Bernoulli}{
El Teorema de Bernoulli es un principio fundamental en la dinámica de fluidos que establece una relación entre la velocidad de un fluido, su presión y su energía potencial. Este teorema se aplica a fluidos incompresibles y en flujo estacionario, y es utilizado en diversas aplicaciones ingenieriles como la aerodinámica, la hidráulica y la ingeniería de tuberías.

El teorema se basa en la conservación de la energía y puede expresarse matemáticamente como:

\begin{equation}
P + \frac{1}{2}\rho v^2 + \rho gh = \text{constante}
\end{equation}

Donde:
\begin{itemize}
    \item $P$: Presión del fluido [Pa]
    \item $\rho$: Densidad del fluido [kg/m^3]
    \item $v$: Velocidad del fluido [m/s]
    \item $g$: Aceleración de la gravedad [9.81 m/s^2]
    \item $h$: Altura relativa en el campo gravitatorio [m]
\end{itemize}

Esta fórmula indica que, en un flujo de fluido ideal, la suma de la presión, la energía cinética por unidad de volumen y la energía potencial por unidad de volumen es constante a lo largo de una línea de corriente.

Una forma más común y práctica de usar el Teorema de Bernoulli es la siguiente, especialmente en problemas de fluidos en movimiento a través de tuberías de diferentes secciones:

\begin{equation}
P_1 + \frac{1}{2}\rho v_1^2 + \rho gh_1 = P_2 + \frac{1}{2}\rho v_2^2 + \rho gh_2
\end{equation}

Donde los subíndices 1 y 2 indican dos puntos diferentes a lo largo de la línea de corriente. Esta ecuación es particularmente útil para resolver problemas en los cuales se conocen las condiciones en un punto y se desea determinar las condiciones en otro.

Otro aspecto importante del Teorema de Bernoulli es su relación con la ecuación de continuidad, que establece que la masa de flujo de un fluido debe conservarse a lo largo de una tubería:

\begin{equation}
A_1 v_1 = A_2 v_2
\end{equation}

Donde:
\begin{itemize}
    \item $A_1$ y $A_2$: Áreas de las secciones transversales en dos puntos de la tubería [m^2]
    \item $v_1$ y $v_2$: Velocidades del fluido en esos puntos [m/s]
\end{itemize}

Este principio de continuidad es esencial para comprender cómo los cambios en la geometría de una tubería afectan la velocidad y la presión del fluido.

}

\actividad{Responder el siguiente cuestionario}
\begin{enumerate}
\item ¿Qué establece el Teorema de Bernoulli?
\item ¿Cuáles son las tres formas de energía que se suman en el Teorema de Bernoulli?
\item ¿En qué tipo de fluidos y condiciones es aplicable el Teorema de Bernoulli?
\item ¿Cómo se relaciona la velocidad de un fluido con su presión según el Teorema de Bernoulli?
\item ¿Qué representa la ecuación de continuidad en el contexto de la dinámica de fluidos?
\item ¿Qué sucede con la velocidad de un fluido cuando el área de la sección transversal de la tubería se reduce?
\item ¿Cómo se puede aplicar el Teorema de Bernoulli en la práctica para medir la velocidad de un fluido?
\item ¿Qué diferencias existen entre un flujo laminar y un flujo turbulento?
\item ¿Cómo afecta la altura en el campo gravitatorio a la presión de un fluido en reposo?
\item ¿Qué aplicaciones prácticas tiene el Teorema de Bernoulli en la ingeniería civil y mecánica?

\end{enumerate}

\actividad{Resolver los siguientes problemas}
\begin{enumerate}
\item En una tubería horizontal, el agua fluye con una velocidad de $3 \, \text{m/s}$ y una presión de $20000 \, \text{Pa}$. Si la velocidad en otro punto de la tubería es de $6 \, \text{m/s}$, ¿Cuál es la presión en ese punto? (Suponga $\rho = 1000 \, \text{kg/m}^3$ y $g = 9.81 \, \text{m/s}^2$)
\item El agua fluye a través de una tubería con un diámetro de $0.1 \, \text{m}$ a una velocidad de $2 \, \text{m/s}$. Si la tubería se estrecha a un diámetro de $0.05 \, \text{m}$, ¿cuál es la velocidad del agua en la sección estrecha?
\item Un fluido de densidad $850 \, \text{kg/m}^3$ fluye desde una altura de $10 \, \text{m}$ hasta el nivel del suelo. Si la velocidad del fluido en la parte superior es cero, ¿cuál es la velocidad del fluido al llegar al suelo?
\item En un sistema de tuberías, el agua tiene una presión de $50000 \, \text{Pa}$ y una velocidad de $1 \, \text{m/s}$ en un punto A. En otro punto B, la presión es de $40000 \, \text{Pa}$. ¿Cuál es la velocidad del agua en el punto B?
\item Un tanque de agua abierto tiene un orificio a $5 \, \text{m}$ por debajo de la superficie del agua. ¿Con qué velocidad sale el agua por el orificio? (Use $\rho = 1000 \, \text{kg/m}^3$ y $g = 9.81 \, \text{m/s}^2$)
\item El aire fluye a través de un tubo horizontal con un diámetro de $0.2 \, \text{m}$ a una velocidad de $5 \, \text{m/s}$. Si la densidad del aire es $1.2 \, \text{kg/m}^3$, ¿cuál es la presión en un punto donde el diámetro del tubo es de $0.1 \, \text{m}$ y la velocidad es de $10 \, \text{m/s}$?
\item Un río tiene una sección transversal de $30 \, \text{m}^2$ y fluye a una velocidad de $2 \, \text{m/s}$. Si el río se estrecha a una sección de $15 \, \text{m}^2$, ¿cuál es la nueva velocidad del flujo?
\item En un laboratorio, un líquido de densidad $950 \, \text{kg/m}^3$ fluye a través de un tubo con una presión de $30000 \, \text{Pa}$ y una velocidad de $4 \, \text{m/s}$. ¿Cuál es la presión en otro punto donde la velocidad del líquido es $2 \, \text{m/s}$?
\item Un avión vuela a una velocidad de $200 \, \text{m/s}$ a nivel del mar, donde la presión atmosférica es $101325 \, \text{Pa}$ y la densidad del aire es $1.225 \, \text{kg/m}^3$. ¿Cuál es la presión en un punto de la superficie del ala donde la velocidad del aire es $150 \, \text{m/s}$?
\item En una tubería vertical, el agua fluye desde una altura de $20 \, \text{m}$ hasta una altura de $5 \, \text{m}$. Si la velocidad del agua en la parte superior es $1 \, \text{m/s}$, ¿cuál es la velocidad en la parte inferior? (Suponga $\rho = 1000 \, \text{kg/m}^3$ y $g = 9.81 \, \text{m/s}^2$)

\end{enumerate}

\end{document}
